(1) 半径 \(s\) 的质量环 \(dM=2\pi\Sigma s\,ds\)，故 \(\Phi=-2\pi G\Sigma\int_0^a s\,ds/\sqrt{s^2+z^2}=-2\pi G\Sigma(\sqrt{a^2+z^2}-z)\)，无穷远极限为零。
(2) \(g_z=-2\pi G\Sigma[1-z/\sqrt{a^2+z^2}]\)，\(dg_z/dz=2\pi G\Sigma a^2/(a^2+z^2)^{3/2}\)。\(z\ll a\) 时 \(g_z\simeq-2\pi G\Sigma(1-z/a)\)；\(z\gg a\) 时 \(g_z\simeq-\pi G\Sigma a^2/z^2\)。
(3) 能量守恒 \(\frac12mv(z)^2+m\Phi(z)=\frac12mv_0^2+m\Phi(0)\)。恰好逃逸要求 \(v_{\rm esc}^2=-2\Phi(0)=4\pi G\Sigma a\)；任意高度 \(v(z)^2=v_0^2+2[\Phi(0)-\Phi(z)]\)，右端为零处即转向点。
(4) \(M=\pi a^2\Sigma\)，远场恰为 \(g_z=-GM/z^2+O(z^{-4})\)。两轨道轴向分离满足 \(\delta\ddot z=(dg_z/dz)\delta z\)，为潮汐拉伸；跨越理想面质量时法向场跳变 \(4\pi G\Sigma\)，与Gauss定律一致。
